% The open problems of Grothendieck's 1973 Buffalo lectures, and what became of them — standalone sub-article.
% Sources: the recordings (« Erwin Schulhoff » YouTube channel, a pseudonym) and the edition's index of results;
% quotations from the recordings are kept to a few words. pdfLaTeX. Compile: tectonic -X compile article.tex.
\documentclass[11pt,a4paper]{amsart}
\usepackage[T1]{fontenc}
\usepackage[utf8]{inputenc}
\usepackage{lmodern,microtype,amssymb,mathtools,booktabs,tabularx,array}
\usepackage[margin=2.7cm]{geometry}
\usepackage[hidelinks]{hyperref}

\newtheorem{problem}{Problem}
\newtheorem{theorem}{Theorem}[section]
\newtheorem{proposition}[theorem]{Proposition}
\newtheorem{lemma}[theorem]{Lemma}
\newtheorem{corollary}[theorem]{Corollary}
\theoremstyle{definition}
\newtheorem{definition}[theorem]{Definition}
\theoremstyle{remark}
\newtheorem*{status}{Status}
\newtheorem{remark}[theorem]{Remark}
\newcommand{\q}[1]{\guillemotleft\,#1\,\guillemotright}
\newcommand{\ts}[3]{\href{https://www.youtube.com/watch?v=#1&t=#2}{#3}}
\newcommand{\Set}{\mathbf{Set}}
\newcommand{\Cat}{\mathbf{Cat}}
\newcommand{\own}{\textsuperscript{(e)}}
\newcolumntype{Y}{>{\raggedright\arraybackslash}X}

\title[The open problems of the Buffalo lectures]{The open problems of Grothendieck's 1973 Buffalo lectures,\\ and what became of them}
\author{Michel Hua}
\address{Commutator, Paris, France}
\email{michel@commutator.io}
\thanks{ORCID 0009-0002-8304-0683.}
\date{7 November 2026}
\subjclass[2020]{01A60; 18F10, 18C35, 14A15, 14L15}
\keywords{Grothendieck, Buffalo 1973, topoi, accessible categories, algebraic theories, fpqc descent, group schemes, open problems}

\begin{document}
\begin{abstract}
In the spring and summer of 1973 Grothendieck gave three courses at the State University of New York at Buffalo, on topos theory, on algebraic geometry and on group schemes, whose recordings are online. In them he asks some twenty questions he cannot answer, states conjectures, and hands problems to the audience. We list them from the recordings, with the second at which each is asked, and say what the literature has made of each. Most are settled. Several went against what he expected: his characterisation of presheaf topoi, whose two conditions instead characterise the essential subtopoi of presheaf topoi, the non-existence of a right adjoint $f^!$, the colimits of accessible categories, the descent of immersions that are not quasi-compact, and his doubt that quasi-coherent sheaves can be pulled back fully faithfully along universal effective epimorphisms. Perrin proved his main conjecture on group schemes in 1976, and the effectivity of finite equivalence relations had been answered by Venken in 1971, two years before he asked it. Some answers are ours; their proofs are given in an appendix and in companion articles. What still looks open concerns large classes of limits in his \q{problem three}, one question that depends on set theory, and descent along non-quasi-compact universal effective epimorphisms in ZFC.
\end{abstract}
\maketitle
\enlargethispage{2pt}

\section{The lectures and the method}

The recordings are on a YouTube channel under the pseudonym \q{Erwin Schulhoff}: four videos for the topos course (April--July 1973), four for algebraic geometry (May--July) and two for group schemes (June--July), the last two also carrying some sessions of the algebraic geometry course. The edition has made an index of results for each course, with the second at which each statement is made, and reconstructed the blackboards \cite{repo}. For this note we searched the transcripts for the places where he says that something is not known, that he cannot prove it, that he expects or suspects it, or where he calls it a conjecture or a problem; we kept the ones he frames as questions, and dropped asides and routine omissions. Each problem below gives the video and the time (a link to that second), how he frames it, and its status, with the references we checked. Where we give an argument of our own, not found in the literature, it is marked \own.

The transcripts are first passes from captions and Whisper that no one has checked, and a few exchanges with the room are hard to attribute. He also corrected himself: on 10 July he retracted several claims of 2 July on algebraic theories (the transitivity of strict cogeneration, a two-step construction, a criterion by strict cogenerators). We report the question as finally posed.

\section{Topoi}

Video 2 of the topos course is \texttt{PqCQM3Pcl6Q} unless stated.

\begin{problem}[\ts{PqCQM3Pcl6Q}{1h26m22s}{1:26:22} and \ts{PqCQM3Pcl6Q}{7h31m23s}{7:31:23}]\label{p:presheaf}
Are the presheaf topoi $\widehat C$ exactly the topoi that share all the exactness properties of $\Set$, including those involving infinite limits: arbitrary sums and products distributive, and arbitrary products of epimorphisms epimorphic?
\end{problem}
He suspects it, says it should not be too hard, and returns to it later as a characterisation of what he calls essential topoi, at least when they have generators.

\begin{status}
False, and the class his two conditions define can be described exactly\own. Write (A) for \q{arbitrary products of epimorphisms are epimorphisms} and (B) for \q{arbitrary products distribute over arbitrary sums}. For a Grothendieck topos $E$, (A) and (B) together are equivalent to each of: $E$ is an essential subtopos of a presheaf topos; for a small generating subcategory with its induced topology, covering sieves are closed under arbitrary intersections; $E$ is totally distributive in the sense of Rosebrugh and Wood. The last equivalence is Lucyshyn-Wright's theorem \cite[1.5.5]{lw}; the others are proved in Appendix~\ref{app:presheaf}, using the axiom of choice. Condition (B) alone says that $E$ is locally connected, and the two conditions are independent: sheaves on the Cantor set satisfy (A) and not (B), sheaves on $\mathbf R$ satisfy (B) and not (A). A topos with both and not a presheaf topos is the topos of sheaves on $[0,1)$ whose open sets are the intervals $[0,a)$: its frame is the complete chain $[0,1]$, which is completely distributive, so (A) and (B) hold, but it has no nonzero projective object, while a presheaf topos is generated by its projective representables; the presheaf topoi are exactly the topoi with (B) and enough projectives (Proposition~\ref{a:proj}). In the localic case, $\mathrm{Sh}(L)$ satisfies (A) and (B) exactly when $L$ is completely distributive, and is a presheaf topos exactly when $L$ is join-generated by completely join-prime elements. His \q{essential} is not the essential geometric morphism of later usage; that the class he named turns out to be that of the essential subtopoi of presheaf topoi looks like a coincidence of words. We do not claim that these equivalences are new, only that we have not found them stated.
\end{status}

\begin{problem}[\ts{PqCQM3Pcl6Q}{4h5m50s}{4:05:50}]\label{p:fshriek}
For a morphism of topoi $f$, can $f_*$ have a right adjoint $f^!$ on sheaves of sets without $f$ being an equivalence?
\end{problem}
He has \q{a feeling} that it cannot, and concludes that the notation $f^!$ has no reason to exist for topoi.

\begin{status}
False, even for an open dense immersion, the case he was testing. For the open point $j$ of the Sierpi\'nski space $S$, whose sheaves are the maps $A\to B$ (sections over $S$, then over the open point), $j_*$ is the diagonal $A\mapsto(A\to A)$, which is the constant-sheaf functor of $S$, and its right adjoint is the global-sections functor $\Gamma$\own. More generally, for a small category $C$ with a terminal object the global-sections functor $\Gamma\colon \widehat C\to\Set$, the direct image of $\widehat C\to\Set$, has a right adjoint: $\widehat C$ is a local topos, and such morphisms are the local maps studied by Johnstone and Moerdijk \cite{jm}.
\end{status}

\begin{problem}[\ts{PqCQM3Pcl6Q}{5h18m59s}{5:18:59}]\label{p:classifying}
Every type of algebraic structure, in a broad sense, should have a classifying topos with a universal model, added axioms should correspond to subtopoi, and the homotopy type of the classifying topos should be a classifying space of the structure.
\end{problem}
He states it as a programme and says he has checked it on a dozen or two examples.

\begin{status}
Carried out for geometric theories: every geometric theory has a classifying topos with a generic model (Joyal--Reyes; Makkai--Reyes \cite{mr}; Hakim for rings and local rings \cite{hakim}), quotient theories correspond to subtopoi (Caramello \cite{caramello}), and for groups and categories the classifying topos has the weak homotopy type of the classifying space (Moerdijk \cite{moerdijk}). The restriction is that \q{structure} means geometric theory.
\end{status}

\begin{problem}[\ts{PqCQM3Pcl6Q}{8h38m47s}{8:38:47}, \ts{PqCQM3Pcl6Q}{8h46m26s}{8:46:26}]\label{p:exactness}
Can every exactness property be written as \q{every $D$-functor $I\to C$ factors through a given $D$-functor $I\to I'$}? Is \q{every semigroup object is a monoid} an exactness property?
\end{problem}
This comes from a discussion with the room after the session, and who says what is not always clear.

\begin{status}
Formalised by Jacqmin and Janelidze as exactness sequents of sketches, which cover what they call first-order exactness properties, among them properties of internal structures such as an internal monoid being a group, and not higher-order ones such as having enough projectives \cite{jacqmin}. \q{Every semigroup object is a monoid} holds only in the trivial category: the left-zero multiplication $X\times X\to X$ is associative, and a unit $e\colon 1\to X$ for it forces $\mathrm{id}_X=e\circ{!}$, so $X\cong 1$; since every object carries this multiplication, the category is trivial\own.
\end{status}

Two further questions are answered on the spot. Asked for a criterion characterising sober spaces (\ts{PqCQM3Pcl6Q}{1h55m37s}{1:55:37}), he says he knows none that can be read off the topos, since a space and its sobrification have the same sheaves; the intrinsic criteria are standard \cite{stone}. And the universal problem for torsors under a group $G$ of a topos $S$ (\texttt{ORG14LNb5gE}, \ts{ORG14LNb5gE}{3h15m58s}{3:15:58}) he solves at once with the topos of $G$-objects of $S$, as in Giraud \cite{giraud}.

\section{Algebraic theories relative to a class of limits}

For a class $D$ of limits and colimits, he considers the 2-category $\Cat_D$ of categories with $D$-limits and $D$-functors, and calls a theory algebraic if it is obtained from a set $I$ of base objects by stepwise rules: the free $D$-envelope, adjoining arrows, imposing axioms.

\begin{problem}[\ts{PqCQM3Pcl6Q}{8h34m57s}{8:34:57}]\label{p:representable}
Are the algebraic theories relative to $D$ exactly the theories $C\mapsto \operatorname{Hom}_D(R,C)$ 2-represented by objects $R$ of $\Cat_D$?
\end{problem}
\begin{status}
Settled: yes for small $R$, which is his own argument and the modern presentation of small categories with limits by generators and relations (sketches, essentially algebraic theories \cite[ch.~3]{ar}); no in general. For $D$ the small limits, $R=\Set$ is representable, with $\Set^{\mathrm{op}}$ as its category of models in $\Set$, which is not accessible, while stepwise theories have locally presentable categories of models; he gives this example himself on 10 July. For $D$ all small limits the algebraic $R$ are, up to duality, the opposites of locally presentable categories, the form Gabriel--Ulmer duality gives \cite{gu}.
\end{status}

\begin{problem}[\ts{PqCQM3Pcl6Q}{10h3m6s}{10:03:06}]\label{p:envelope}
For which $D$ are the envelope problem (the free $D$-completion of a category) and the localisation problem in $\Cat_D$ solvable? He knows them for finite limits and for topoi, and not for additive, abelian or Karoubian categories.
\end{problem}
\begin{status}
Settled for every small class $D$ by Kelly's transfinite constructions and two-dimensional monad theory \cite{kelly, bkp}, with explicit answers: the idempotent completion for Karoubian categories, the biproduct completion of the free preadditive category for additive ones, the free abelian category of Freyd and Adelman \cite{freyd, adelman}, the exact completion of Carboni and Celia Magno \cite{ccm}. For a large class, such as all small limits, the answers depend on size, as he noted on 10 July: localisations of large categories need not be locally small.
\end{status}

\begin{problem}[\texttt{4gynb9uTEE4}, \ts{4gynb9uTEE4}{2h3m54s}{2:03:54} and \ts{4gynb9uTEE4}{2h10m23s}{2:10:23}]\label{p:three}
\q{Problem three}: is $I\to R$ basic if and only if every strictly full subcategory of $R$ stable under $D$-limits and containing the image of $I$ is $R$ itself? And, conjecturally, can $I$ be enlarged so that no arrows remain to be adjoined, only axioms?
\end{problem}
\begin{status}
Settled, except for mixed large classes $D$\own; the proofs are in the companion article \cite{three}. Write (G) for the condition of the problem; it holds exactly when the transfinite closure of $I$ under $D$-limits and $D$-colimits exhausts $R$, after $\rho$ steps when the shapes of $D$ have fewer than $\rho$ arrows \cite[\S3.5]{kellybook}. Assuming, as he does, that envelopes, localisations and the other constructions exist in $\Cat_D$, basic implies (G) for every $D$, and $I\to R$ is basic if and only if (G) holds and $R$ is algebraic over some base \cite{three}. For a small class $D$ every $R$ satisfying (G) is essentially small, hence algebraic, so the equivalence holds with the base fixed; it also holds for Grothendieck topoi. For $D$ all small limits it fails: every set is an equaliser of two maps between powers of $2$, so $\{2\}\to\Set$ satisfies (G) after two steps, but $\Set$ is not algebraic, since $2$ is not presentable in $\Set^{\mathrm{op}}$. The statements about sets are checked in Lean, with the standard axioms only (\texttt{lean/\allowbreak Buffalo/\allowbreak ProblemThree.lean} in \cite{repo}): every set is in bijection with that equaliser, the saturation of $\{2\}$ under limits has every set at its second stage and no three-element set at its first, and $2$ is not $\kappa$-presentable in $\Set^{\mathrm{op}}$ for any regular $\kappa$; that $\Set$ is not algebraic, which uses Kelly's sketches, is not formalised. In an algebraic category the saturation stops after a small number of steps, so an example of Adámek and Rosický \cite[Example~6.38]{ar}, a colimit closure that needs a proper class of steps, gives a second counterexample, for all small colimits, and shows that the number of steps is not bounded in general \cite{three}. The second question has a positive answer in the locally presentable case by Gabriel--Ulmer: taking for $I$ the opposite of the $\lambda$-presentable objects, the category of models is the category of $\lambda$-continuous functors on $I$, so that only axioms remain to be imposed \cite{gu}. What remains open is the converse for mixed large classes, such as finite limits with all colimits but no exactness condition. Two related statements of the course on topoi (\texttt{ORG14LNb5gE}), that every object of a later stage is covered by images of earlier ones and that the finite products of a basic $I$ generate, are false as stated: adjoining an arrow $1\to1\sqcup1$ to $\Set$ gives $\Set\times\Set$, in which $(1,0)$ is not covered by the diagonal \cite{three}.
\end{status}

Around these he asks whether faithfulness of $\operatorname{Hom}_D(R,C)\to C^I$ for all $C$ suffices for algebraicity (\ts{PqCQM3Pcl6Q}{9h40m36s}{9:40:36}, \ts{PqCQM3Pcl6Q}{10h0m19s}{10:00:19}). It does not, as he reports himself later, citing an example he attributes to Isbell (the name is uncertain in the recording; \ts{4gynb9uTEE4}{2h3m54s}{2:03:54}). In the session itself he and the room arrive at lattices, and a large ordinal; $\Set$ with $I=\{2\}$ is another counterexample, faithfulness holding because every set is a limit of powers of $2$\own. His hope of a bound on the number of limit operations (\ts{PqCQM3Pcl6Q}{8h8m14s}{8:08:14}) he retracted himself on 10 July.

\section{Accessible categories}

\begin{problem}[\texttt{4gynb9uTEE4}, \ts{4gynb9uTEE4}{1h31m3s}{1:31:03}]\label{p:colim}
Is a small 2-colimit in $\Cat$ of accessible categories accessible?
\end{problem}
\begin{status}
No in general\own. Let $B$ be the free category on arrows $a\colon u\to v$, $b\colon v\to w$, $h\colon u\to w$, map it to the chain $0<1<2$ and to the free monoid on one generator $t$, sending $a,b,h$ to $t$. The pushout, strict or pseudo, is the one-object category with endomorphisms $\{1,t\}$ and $t=t^2$, since $h=ba$ in the chain; the idempotent $t$ does not split, while accessible categories split idempotents. The three categories are small and idempotent-complete, hence accessible, and the functors are accessible. He expected the answer yes, while noting that a sum of two categories with colimits has none. Directed colimits of accessible categories along accessible full embeddings are accessible (Paré--Rosický \cite{pr}). The limit side, which he recalls from SGA~4, is right in modern form \cite[ch.~5]{mp}.
\end{status}

\begin{problem}[\texttt{4gynb9uTEE4}, \ts{4gynb9uTEE4}{1h51m31s}{1:51:31}]\label{p:reflective}
Is a full subcategory of an accessible category, stable under limits and with a left adjoint inclusion, accessible?
\end{problem}
\begin{status}
Yes under an extra hypothesis; in ZFC alone we have not found it settled. When the inclusion is also accessible, the condition he proposed adding, the answer is yes in ZFC: for a locally presentable ambient category this is \cite[2.78]{ar}. Under Vop\v{e}nka's principle the answer is yes for a locally presentable ambient category \cite[ch.~6]{ar}. Under the weak Vop\v{e}nka principle every limit-closed full subcategory of a locally presentable category is reflective, and this fails in the category of graphs when the weak principle fails \cite{art}, \cite[6.22]{ar}; Wilson showed that the weak principle does not imply the full one \cite{wilson}. A reflective, limit-closed, non-accessible subcategory in a model of the weak principle without the full one would show that ZFC does not decide the question; we have not found one in print.
\end{status}

\begin{problem}[\texttt{4gynb9uTEE4}, \ts{4gynb9uTEE4}{4h48m17s}{4:48:17}]\label{p:dual}
If $E$ and $E^{\mathrm{op}}$ are both accessible, is $E$ equivalent to the point?
\end{problem}
\begin{status}
No: the two-element chain is a counterexample, being, like its opposite, small and idempotent-complete. With small colimits in $E$, as he assumes for his main theorem, $E$ and $E^{\mathrm{op}}$ are both locally presentable, and then $E$ is a preorder, hence a complete lattice (Gabriel--Ulmer \cite[Satz~7.13]{gu}).
\end{status}

\section{Algebraic geometry and group schemes}

\begin{problem}[\texttt{vuys5SNs2eQ}, \ts{vuys5SNs2eQ}{3h31m11s}{3:31:11}; see also \ts{vuys5SNs2eQ}{2h59m31s}{2:59:31}]\label{p:immersion}
Do immersions that are not quasi-compact descend along a faithfully flat quasi-compact morphism? At 2:59:31: is descent of subschemes exact without retrocompactness? (\q{you asked for an open problem}).
\end{problem}
He believes the first but cannot prove it; his attempt by a filtered union of quasi-compact opens fails on his own example, the closed points of the line.

\begin{status}
Both are false. Open, closed and quasi-compact immersions descend along fpqc coverings, and immersions along fppf ones \cite[02L3, 02L6, 02L8, 02YM]{stacks}; on the underlying sets, locally closed and retrocompact subsets descend along fpqc morphisms, and locally closed ones along fppf morphisms \cite[Thms~4.1, 4.2]{rydh}. The disjoint union of the $\operatorname{Spec}\mathbf F_p$ mapped to $\operatorname{Spec}\mathbf Z$ is not an immersion, but becomes one after the faithfully flat base change $\operatorname{Spec}\prod_p\mathbf Z_{(p)}\to\operatorname{Spec}\mathbf Z$; the same subscheme answers the second question. The example, base change included ($\prod_q\mathbf Z_{(q)}/p\cong\mathbf F_p$), is in a comment of August 2026 on Stacks tag 02YL, which we checked; it is not refereed, and we found no earlier published counterexample. The subscheme is the one he used in class, the closed points, taken after base change.
\end{status}

\begin{problem}[\texttt{3ULFIl\_2Wis}, \ts{3ULFIl_2Wis}{5h4m48s}{5:04:48}]\label{p:perrin}
\q{The main conjecture}: a connected separable group scheme over a field is an extension of a group of finite type by an affine group, the quotient map faithfully flat and quasi-compact; equivalently a filtered limit of groups of finite type with affine transition maps.
\end{problem}
\begin{status}
Proved by Perrin \cite{perrin}, a summary of his Orsay work: every quasi-compact group scheme over a field is a filtered limit of algebraic quotients with affine kernels, with transition maps affine from some index on; he notes that the statement had been conjectured in EGA~IV. A connected group scheme over a field is quasi-compact, as Grothendieck recalls just after the statement. Perrin's counterexample with abelian varieties of growing dimension is the phenomenon Grothendieck wanted to exclude at \ts{3ULFIl_2Wis}{5h11m21s}{5:11:21}, the dimensions of the abelian quotients being bounded.
\end{status}

Asked for a characterisation of schemes among algebraic spaces (\ts{3ULFIl_2Wis}{6h39m12s}{6:39:12}, a session of the algebraic geometry course recorded on this video), he answers that it would solve the representability problem, and no; there are sufficient criteria and partial results \cite[03XX, 0ADD, 03JI, 0D2W]{stacks}. Smaller questions are settled. In the algebraic geometry course: whether a Noetherian spectrum, the ascending chain condition on radical ideals, is the same as the ascending chain condition on primes (\texttt{51wmmpHand0}, \ts{51wmmpHand0}{7h42m27s}{7:42:27}; no: an infinite product of fields has only maximal primes and an infinite, non-Noetherian spectrum\own; see Ohm and Pendleton for rings with Noetherian spectrum \cite{ohm}); and whether a scheme can have no closed point (\texttt{Oqhr-Gt66es}, \ts{Oqhr-Gt66es}{0h10m49s}{0:10:49}; yes \cite[01IY]{stacks}). There he also states that the local schemes are those with a unique closed point. That holds for quasi-compact schemes, but not in general: $W\sqcup\operatorname{Spec}k$, with $W$ a scheme without closed point, has a unique closed point and is not local\own. In the group-schemes course: whether a monoid scheme with a closed unit section is separated (\texttt{3ULFIl\_2Wis}, \ts{3ULFIl_2Wis}{3h56m26s}{3:56:26}; no: adjoin a unit, as an extra rational point, to a non-separated scheme with a rational point $c$ and the constant multiplication $xy=c$\own). 

\begin{problem}[\texttt{3ULFIl\_2Wis}, \ts{3ULFIl_2Wis}{6h26m20s}{6:26:20}--\ts{3ULFIl_2Wis}{6h27m24s}{6:27:24}]\label{p:qcoh}
Along a morphism of universal one-descent, that is a universal effective epimorphism of schemes, do quasi-coherent sheaves descend? Is pulling them back at least fully faithful?
\end{problem}
A listener asks why the topology is not defined by all such morphisms; he answers that it would be the finest one for which representable functors are sheaves, but that these morphisms are hard to handle, that one does not know whether quasi-coherent sheaves descend along them or whether pulling back is even faithful, and that he doubts this can be proved. This is an algebraic geometry session recorded on the group-schemes video.

\begin{status}
Mostly settled\own, and partly against his doubt; the proofs are in a companion article \cite{qcoh}. Pulling back quasi-coherent sheaves along any universal effective epimorphism is fully faithful, with no quasi-compactness: test against the square-zero extension $\mathcal O_Y\oplus F$ and the affine line. Along a quasi-compact one, descent is effective; this is due to Mesablishvili \cite{mesablishvili2004,mesablishvili-acs}, and the companion article gives an independent proof through a finite affine cover, which reduces to the affine case, where universal effective epimorphisms are the pure ring maps and descent holds exactly along pure maps \cite[0EUA, 08XA]{stacks}, \cite{mesablishvili}. Without quasi-compactness, effectivity can fail: if some nonprincipal ultrafilter $U$ on the primes is generated by $\aleph_1$ sets, for instance under the continuum hypothesis, the family $\{\operatorname{Spec}\mathbf Z_A\to\operatorname{Spec}\mathbf Z\}_{A\in U}$ of the Stacks Project \cite[0EUG]{stacks} is a universal effective epimorphism along which some descent datum, an extension of $\mathcal O$ by $\bigoplus_p\mathbf Z/p$, is not effective; the obstruction is a clopen set of $\beta\mathbf N\setminus\mathbf N$ minus the point $U$ that does not extend, which exists at points of character $\aleph_1$ by the butterfly lemma of Fine and Gillman \cite{pitz}. In ZFC alone the question is open: in the model of van Douwen, Kunen and van Mill the obstruction used here vanishes. The affine core of full faithfulness is checked in Lean, with the standard axioms only (\texttt{lean/\allowbreak Buffalo/\allowbreak QcohDescent.lean} in \cite{repo}): if $N\to B\otimes_AN\rightrightarrows B\otimes_AB\otimes_AN$ is an equaliser, base change $\operatorname{Hom}_A(M,N)\to\operatorname{Hom}_{\mathrm{Des}}(B\otimes_AM,\,B\otimes_AN)$ is bijective for every $A$-module $M$; conversely, bijectivity for $M=A$ gives the equaliser condition for $N$, and that condition for all modules is equivalent to $A\to B$ being pure \cite[08WE]{stacks}. The step from universal effective epimorphisms to the equaliser condition, effectivity, and the non-affine case are not formalised.
\end{status}

\section{Effectivity}\label{sec:eff}

At the end of the algebraic geometry session of 11 July, recorded on the second group-schemes video (\texttt{3ULFIl\_2Wis}, \ts{3ULFIl_2Wis}{6h59m50s}{6:59:50}), asked for open problems, he gives the effectivity of finite equivalence relations on a finite scheme over a field, posed in FGA in 1961, and says he suspects a negative answer. It had been answered negatively, over every field, by Venken in 1971: Corollary~3.1 of a note on the descent of flat modules gives a non-effective equivalence relation on the spectrum of a local Artin algebra of dimension $33$ \cite{venken}. A companion note checks that example by machine and gives a second one, of length $15$ ($16$ in characteristic~$2$), by restricting Raicu's relation on the plane \cite{effective}.

\section{Summary}

\begin{table}[ht]
\small
\begin{tabularx}{\linewidth}{@{}l>{\raggedright\arraybackslash}p{5.3cm}Y@{}}
\toprule
Problem & His stance & Status \\
\midrule
\ref{p:presheaf} presheaf topoi & suspects yes & false: Scott sheaves on $[0,1]$\own; essential subtopoi of $\widehat C$ \cite{lw} \\
\ref{p:fshriek} $f^!$ & feels no & false; local morphisms \cite{jm} \\
\ref{p:classifying} classifying topoi & programme & done for geometric theories \\
\ref{p:exactness} exactness properties & open in discussion & formalised \cite{jacqmin}; semigroup case trivial\own \\
\ref{p:representable} algebraic $=$ representable & aim & yes for small $R$, no in general \\
\ref{p:envelope} envelopes, localisations & known in two cases & small $D$: yes \cite{kelly,bkp} \\
\ref{p:three} problem three & neither direction proved & $\Rightarrow$ yes; $\Leftrightarrow$ for small $D$; fails for all limits\own \\
\ref{p:colim} colimits of accessible & expects yes & no\own; directed, full embeddings: yes \cite{pr} \\
\ref{p:reflective} reflective subcategories & asks for a counterexample & yes if the inclusion is accessible, or under Vop\v{e}nka \cite{ar}; ZFC alone open to us \\
\ref{p:dual} $E$, $E^{\mathrm{op}}$ accessible & guesses the point & no: the two-element chain \cite{gu} \\
\ref{p:immersion} descent of immersions & believes yes & no (Stacks comment, 2026) \\
\ref{p:perrin} structure of group schemes & main conjecture & Perrin 1976 \cite{perrin} \\
\ref{p:qcoh} descent of quasi-coherent sheaves & doubts it & fully faithful; effective if quasi-compact; can fail under CH\own \\
\S\ref{sec:eff} effectivity & suspects no & no, over every field (Venken 1971) \cite{venken,effective} \\
\bottomrule
\end{tabularx}
\end{table}

The questions that still look open are the converse of problem three for mixed large classes $D$ (Problem~\ref{p:three}), the status in ZFC alone of Problem~\ref{p:reflective}, and the effectivity of descent of quasi-coherent sheaves along non-quasi-compact universal effective epimorphisms in ZFC (Problem~\ref{p:qcoh}). Six of his expectations turned out wrong (Problems~\ref{p:presheaf}, \ref{p:fshriek}, \ref{p:colim}, \ref{p:dual}, \ref{p:immersion}, and his doubt in Problem~\ref{p:qcoh} about full faithfulness); his programme of classifying topoi, his main conjecture and his doubt about effectivity were borne out.

\subsection*{Use of language models}
The questions were located in the transcripts, and the literature searched, with language models (Claude Opus~5.5, Anthropic), which also wrote a first draft of this text and the arguments marked \own, under the author's direction. The references were opened on their pages; the arguments marked \own\ have not been checked by a specialist.

\appendix

\section{Presheaf topoi and the conditions (A) and (B)}\label{app:presheaf}

This appendix proves the claims made under Problem~\ref{p:presheaf}. Throughout, $E$ is a Grothendieck topos, the metatheory is ZFC, and the axiom of choice is used freely. The two conditions are:
\begin{itemize}
\item[(A)] for every small family of epimorphisms $f_i\colon X_i\to Y_i$, the product $\prod_i f_i\colon\prod_i X_i\to\prod_i Y_i$ is an epimorphism;
\item[(B)] for every small family of small families $(X_{ij})_{i\in I,\,j\in J_i}$, the canonical map
\[\textstyle\coprod_{\varphi\in\prod_i J_i}\prod_i X_{i\varphi(i)}\longrightarrow\prod_i\coprod_{j\in J_i}X_{ij}\]
is an isomorphism.
\end{itemize}
Following Rosebrugh and Wood \cite{rw}, as recalled by Lucyshyn-Wright \cite[\S1]{lw}, a locally small category $E$ is \emph{totally distributive} if there is a chain of adjoint functors $t\dashv c\dashv y\colon E\to\widehat E=[E^{\mathrm{op}},\Set]$ with $y$ the Yoneda embedding, that is, if $c=\operatorname{colim}$ has a left adjoint. An \emph{essential subtopos} of a topos $F$ is a topos $E$ with a geometric inclusion $i\colon E\to F$ whose inverse image $i^*$ has a left adjoint \cite[\S1.5]{lw}. Lucyshyn-Wright's Theorem~1.5.5 \cite{lw} states that the totally distributive categories having a small dense generator are exactly the essential subtopoi of presheaf topoi $\widehat C$, $C$ small; every Grothendieck topos has a small dense generator, so the theorem applies to all of them. By a theorem of Kelly and Lawvere \cite{kl}, recalled in \cite[1.5.6]{lw}, the essential subtopoi of $\widehat C$ correspond bijectively to the idempotent ideals of arrows of $C$.

\subsection{The class defined by (A) and (B)}

\begin{proposition}\label{a:B}
$E$ satisfies (B) if and only if $E$ is locally connected.
\end{proposition}

\begin{proof}
Assume (B). With all $X_{ij}=1$, condition (B) says that the constant-sheaf functor $\Delta\colon\Set\to E$ preserves small products. Being left exact, $\Delta$ then preserves all small limits; it also preserves all colimits, so it is accessible. Since $\Set$ and $E$ are locally presentable, the adjoint functor theorem \cite[1.66]{ar} gives a left adjoint $\pi_0\dashv\Delta$. Over $\Set$ such a left adjoint is automatically $\Set$-indexed: $X\times\Delta I=\coprod_I X$, so $\pi_0(X\times\Delta I)=\pi_0X\times I$. Hence $E$ is locally connected.

Conversely, a locally connected topos has a generating set of connected objects, the components of the objects of a generating set. For a connected object $Q$, the functor $\operatorname{Hom}(Q,-)$ preserves products, and it preserves coproducts: a map $Q\to\coprod_j X_j$ pulls the summands back to a decomposition $Q=\coprod_j Q_j$, and exactly one $Q_j$ is nonzero, since $Q\neq0$ and $Q=Q_{j_0}\sqcup\coprod_{j\neq j_0}Q_j$. So $\operatorname{Hom}(Q,-)$ sends the canonical map of (B) to the canonical bijection
\[\textstyle\coprod_\varphi\prod_i\operatorname{Hom}(Q,X_{i\varphi(i)})\cong\prod_i\coprod_j\operatorname{Hom}(Q,X_{ij})\]
of $\Set$, which uses the axiom of choice. A map of $E$ that is bijective on $\operatorname{Hom}(Q,-)$ for every $Q$ in a generating set is an isomorphism: it is a monomorphism because the generators separate maps, and an epimorphism because every map from a generator to its target lifts.
\end{proof}

\begin{proposition}\label{a:ess}
An essential subtopos of a presheaf topos satisfies (A) and (B).
\end{proposition}

\begin{proof}
Let $i\colon E\hookrightarrow\widehat C$ be the inclusion, with $i_!\dashv i^*\dashv i_*$. Since $i_*$ is fully faithful, so is $i_!$, and $\mathrm{id}\cong i^*i_!$.

(A): each $i_!f_i$ is an epimorphism, as the image of one under a left adjoint. The product $\prod_i i_!f_i$ is an epimorphism of $\widehat C$, since products are computed pointwise and the axiom of choice holds in $\Set$. The functor $i^*$ preserves products and epimorphisms, so $\prod_i f_i\cong i^*\bigl(\prod_i i_!f_i\bigr)$ is an epimorphism.

(B): naturally, and compatibly with the canonical maps,
\[\textstyle\prod_i\coprod_jX_{ij}\cong i^*\bigl(\prod_i\coprod_j i_!X_{ij}\bigr)\cong i^*\bigl(\coprod_\varphi\prod_i i_!X_{i\varphi(i)}\bigr)\cong\coprod_\varphi\prod_iX_{i\varphi(i)}.\qedhere\]
\end{proof}

\begin{proposition}\label{a:td}
A totally distributive Grothendieck topos satisfies (A) and (B).
\end{proposition}

This also follows from Proposition~\ref{a:ess} and Lucyshyn-Wright's theorem; the proof below does not use the theorem.

\begin{proof}
Suppose that $c=\operatorname{colim}\colon\widehat E\to E$ has a left adjoint, so that $c$ preserves all limits that exist in $\widehat E$; products in $\widehat E$ are pointwise.

(B): in $\widehat E$, $\prod_i\coprod_j yX_{ij}\cong\coprod_\varphi\prod_i yX_{i\varphi(i)}=\coprod_\varphi y\bigl(\prod_iX_{i\varphi(i)}\bigr)$, by pointwise distributivity and because $y$ preserves limits. Applying $c$, which preserves products and coproducts and satisfies $c\circ y\cong\mathrm{id}$, gives (B).

(A): let $P_i\subseteq yY_i$ be the sieve generated by $f_i$. In a topos the colimit of a sieve on $Y$ is the union of the images of its members, because the equivalence relation generated by the diagram contains the kernel pair, sieves being closed under precomposition. So $c(P_i)=\operatorname{im}f_i=Y_i$. By the axiom of choice, $\prod_iP_i\subseteq y\bigl(\prod_iY_i\bigr)$ is the sieve generated by $\prod_if_i$. Hence $\operatorname{im}\bigl(\prod_if_i\bigr)=c\bigl(\prod_iP_i\bigr)=\prod_ic(P_i)=\prod_iY_i$.
\end{proof}

\begin{lemma}\label{a:inter}
Let $(C,J)$ be a small site, with associated sheaf functor $a\colon\widehat C\to\mathrm{Sh}(C,J)$. If each $J(c)$ is closed under arbitrary small intersections, then $a$ preserves all small limits; hence it has a left adjoint, and $\mathrm{Sh}(C,J)\hookrightarrow\widehat C$ is an essential subtopos. Conversely, if this inclusion is essential, each $J(c)$ is closed under arbitrary intersections.
\end{lemma}

\begin{proof}
Write $F^+(c)=\operatorname{colim}_{S\in J(c)}\mathrm{Match}(S,F)$; the colimit is directed, since $J(c)$ is closed under finite intersections. Since $\mathrm{Match}(S,-)=\operatorname{Hom}(S,-)$ preserves products, there is a canonical map
\[\textstyle\bigl(\prod_iF_i\bigr)^+(c)=\operatorname{colim}_S\prod_i\mathrm{Match}(S,F_i)\longrightarrow\prod_i\operatorname{colim}_S\mathrm{Match}(S,F_i).\]
It is surjective: represent each $\xi_i$ on some $S_i$, using the axiom of choice to choose the representatives, and restrict them all to $T=\bigcap_iS_i\in J(c)$. It is injective: if two families agree in each factor, they agree on some $S_i$ for each $i$, hence on $\bigcap_iS_i\in J(c)$. So $(-)^+$ preserves products, and so does $a=(-)^{++}$. Being left exact, $a$ preserves all small limits. It is a left adjoint between locally presentable categories, hence accessible, and the adjoint functor theorem \cite[1.66]{ar} gives a left adjoint.

Conversely, suppose that $a$ preserves limits, and let $S_i$ be covering sieves on $c$, so that each $a(S_i)\to a(yc)$ is an isomorphism. Then $a\bigl(\bigcap_iS_i\bigr)$, which is the wide pullback of the $a(S_i)$ over $a(yc)$, maps isomorphically to $a(yc)$, and $\bigcap_iS_i$ is covering.
\end{proof}

\begin{lemma}\label{a:AB-ii}
Let $C\subseteq E$ be a small generating full subcategory with its induced topology, in which a sieve on $c$ covers if and only if its members are jointly epimorphic in $E$. If $E$ satisfies (A) and (B), each $J(c)$ is closed under arbitrary intersections.
\end{lemma}

\begin{proof}
Let $(S_i)_{i\in I}$ be covering sieves on $c$, so that each $e_i\colon\coprod_{(d,s)\in S_i}d\to c$ is an epimorphism.
\begin{enumerate}
\item By (A), $\prod_ie_i$ is an epimorphism onto $c^I$.
\item By (B), its domain is $\coprod_{\varphi\in\prod_iS_i}\prod_id_{\varphi(i)}$, and on the summand $\varphi$ the map is $\prod_is_{\varphi(i)}$.
\item Pull back along the diagonal $c\to c^I$. Epimorphisms and coproducts are stable under pullback in a topos, so $\coprod_\varphi P_\varphi\to c$ is an epimorphism, where $P_\varphi$ is the wide pullback over $c$ of the maps $s_{\varphi(i)}\colon d_{\varphi(i)}\to c$.
\item For each $i$, the map $P_\varphi\to c$ factors through $s_{\varphi(i)}\in S_i$.
\item Cover each $P_\varphi$ by objects of $C$. The composites $e\to P_\varphi\to c$ lie in every $S_i$, hence in $\bigcap_iS_i$, and they are jointly epimorphic. So $\bigcap_iS_i$ covers.\qedhere
\end{enumerate}
\end{proof}

\begin{theorem}\label{a:main}
For a Grothendieck topos $E$ the following are equivalent:
\begin{enumerate}
\item[(i)] $E$ satisfies (A) and (B);
\item[(ii)] for some (equivalently, every) small generating full subcategory $C$ of $E$, the induced topology has each $J(c)$ closed under arbitrary intersections;
\item[(iii)] $E$ is an essential subtopos of a presheaf topos;
\item[(iv)] $E$ is totally distributive.
\end{enumerate}
\end{theorem}

\begin{proof}
(i)$\Rightarrow$(ii) for every $C$ is Lemma~\ref{a:AB-ii}. (ii)$\Rightarrow$(iii): by the comparison lemma $E\simeq\mathrm{Sh}(C,J)$; apply Lemma~\ref{a:inter}. (iii)$\Rightarrow$(i) is Proposition~\ref{a:ess}. (iii)$\Leftrightarrow$(iv) is Lucyshyn-Wright's Theorem~1.5.5 \cite{lw}; the implication (iv)$\Rightarrow$(i) is also Proposition~\ref{a:td}.
\end{proof}

The equivalence of (i), (ii) and (iii) rests only on the axiom of choice and standard facts (the comparison lemma, $a=(-)^{++}$, the adjoint functor theorem for locally presentable categories). We have not found (i)$\Leftrightarrow$(iii) or (i)$\Leftrightarrow$(iv) stated; Lemma~\ref{a:inter} is of the kind Kelly and Lawvere's paper \cite{kl} may contain, and we do not claim novelty.

\subsection{The localic case}\label{app:localic}

Let $L$ be a frame, with the canonical topology, in which covers are joins: a down-set $S$ of $u$ covers $u$ if $\bigvee S=u$. Recall that $L$ is \emph{completely distributive} if $\bigwedge_i\bigvee_jf_{ij}=\bigvee_{\varphi\in\prod_iJ_i}\bigwedge_if_{i\varphi(i)}$ for every family $(f_{ij})$; under the axiom of choice this is the same as constructive complete distributivity, as recalled in \cite[\S1.1]{lw} from Fawcett and Wood \cite{fw}. The inequality $\geq$ always holds.

\begin{lemma}\label{a:cd-cover}
$L$ is completely distributive if and only if, for every $u\in L$, the covering down-sets of $u$ are closed under arbitrary intersections.
\end{lemma}

\begin{proof}
Suppose $L$ completely distributive, and let $S_i$ be down-sets with $u\le\bigvee S_i$ for each $i$. For $\varphi\in\prod_iS_i$, the element $\bigwedge_i\varphi(i)$ lies below $\varphi(i)\in S_i$ for every $i$, so it lies in $\bigcap_iS_i$. Hence $u\le\bigwedge_i\bigvee S_i=\bigvee_\varphi\bigwedge_i\varphi(i)\le\bigvee\bigcap_iS_i$.

Conversely, let $(f_{ij})$ be a family and $u=\bigwedge_i\bigvee_jf_{ij}$. By frame distributivity $u=\bigvee_j(u\wedge f_{ij})$ for each $i$, so the down-set $S_i$ generated by the $u\wedge f_{ij}$, $j\in J_i$, covers $u$. By hypothesis $\bigcap_iS_i$ covers $u$. Each $w\in\bigcap_iS_i$ satisfies, for each $i$, $w\le u\wedge f_{ij}$ for some $j$; choosing such a $j=\varphi(i)$ gives $w\le\bigwedge_if_{i\varphi(i)}$. Hence $u=\bigvee\bigcap_iS_i\le\bigvee_\varphi\bigwedge_if_{i\varphi(i)}$.
\end{proof}

The two directions can also be checked on (A) and (B) directly.

\emph{Complete distributivity implies (A) for $\mathrm{Sh}(L)$.} Let $f_i\colon X_i\to Y_i$ be epimorphisms and $s=(s_i)$ a section of $\prod_iY_i$ over $u$. Each $s_i$ lifts to $X_i$ locally, on a cover $\{u_{ik}\}_k$ of $u$. By complete distributivity $u=\bigwedge_i\bigvee_ku_{ik}=\bigvee_\varphi\bigwedge_iu_{i\varphi(i)}$, and over each $\bigwedge_iu_{i\varphi(i)}$ every $s_i$ lifts, by restriction of its lift over $u_{i\varphi(i)}$, so that $s$ lifts to $\prod_iX_i$. These elements cover $u$, so $\prod_if_i$ is an epimorphism.

\emph{(A) and (B) imply complete distributivity.} Let $u=\bigvee_ku_{ik}$ for each $i$, with $u_{ik}\le u$. Apply (A) to the epimorphisms $\coprod_ky(u_{ik})\to y(u)$, whose product has codomain $y(u)$. The identity section over $u$ lifts locally: $u$ is covered by elements $v$ such that, for each $i$, $v$ has a disjoint decomposition $v=\bigsqcup_kv_{ik}$ with $v_{ik}\le u_{ik}$. Condition (B) applied to the subterminals $y(v_{ik})$ gives $y(v)=\prod_i\coprod_ky(v_{ik})\cong\coprod_\varphi y\bigl(\bigwedge_iv_{i\varphi(i)}\bigr)$, so $v=\bigvee_\varphi\bigwedge_iv_{i\varphi(i)}\le\bigvee_\varphi\bigwedge_iu_{i\varphi(i)}$. Taking the join over the $v$ gives $u\le\bigvee_\varphi\bigwedge_iu_{i\varphi(i)}$. A family $(f_{ij})$ reduces to this case with $u=\bigwedge_i\bigvee_jf_{ij}$ and $u_{ij}=u\wedge f_{ij}$.

The same argument gives the exact content of (A) for $\mathrm{Sh}(L)$: (A) holds if and only if, for every $u$ and every family of covers $\{u_{ik}\}_k$ of $u$, the element $u$ is covered by elements $v$ having, for each $i$, a disjoint decomposition subordinate to the $i$-th cover.

\begin{proposition}\label{a:loc-B}
The following are equivalent: $\mathrm{Sh}(L)$ satisfies (B); $L$ is locally connected; $L$ is \emph{disjointly completely distributive}, that is, $b=\bigvee_\varphi\bigwedge_if_{i\varphi(i)}$ whenever $b=\bigsqcup_jf_{ij}$ is a disjoint decomposition for each $i$. In particular every completely distributive frame is locally connected.
\end{proposition}

\begin{proof}
Locally connected implies (B) by Proposition~\ref{a:B}. (B) implies disjoint complete distributivity by applying (B) to the subterminals: $y(b)=\prod_i\coprod_jy(f_{ij})\cong\coprod_\varphi y\bigl(\bigwedge_if_{i\varphi(i)}\bigr)$.

Suppose $L$ disjointly completely distributive, and fix $b$. Index all disjoint decompositions $D$ of $b$ (they form a set), and for a choice $\varphi$ of a member $\varphi(D)\in D$ for each $D$ put $e_\varphi=\bigwedge_D\varphi(D)$. By hypothesis $b=\bigvee_\varphi e_\varphi$, and $e_\varphi\wedge e_\psi=0$ for $\varphi\neq\psi$, since they differ at some $D$ and distinct members of $D$ are disjoint; so $b=\bigsqcup_\varphi e_\varphi$ is again a disjoint decomposition, the common refinement. Each nonzero $e_\varphi$ is connected: if $e_\varphi=c_1\sqcup c_2$, then $D'=\{c_1,c_2\}\cup\{e_\psi:\psi\neq\varphi\}$ is a disjoint decomposition of $b$, so $e_\varphi\le\varphi(D')$; as $e_\varphi$ is nonzero and disjoint from the $e_\psi$, $\psi\ne\varphi$, the member $\varphi(D')$ is $c_1$ or $c_2$, say $c_1$, and then $c_2=c_2\wedge e_\varphi\le c_2\wedge c_1=0$. So $b$ is a join of connected elements, and $L$ is locally connected. The last statement follows, since complete distributivity contains the disjoint case.
\end{proof}

\begin{corollary}\label{a:loc}
For a locale $L$ the following are equivalent: $\mathrm{Sh}(L)$ satisfies (A) and (B); $L$ is completely distributive; $\mathrm{Sh}(L)\hookrightarrow\mathrm{PSh}(L)$ is essential; $\mathrm{Sh}(L)$ is totally distributive. Moreover $\mathrm{Sh}(L)$ is a presheaf topos if and only if $L$ is the frame of down-sets of a poset, that is, if and only if $L$ is join-generated by completely join-prime elements.
\end{corollary}

\begin{proof}
The subterminals form a small generating full subcategory of $\mathrm{Sh}(L)$, isomorphic to $L$, whose induced topology is the canonical one. Lemma~\ref{a:cd-cover}, Lemma~\ref{a:inter} and Theorem~\ref{a:main} give the first equivalences.

If $\mathrm{Sh}(L)\simeq\widehat C$, the representables are connected projective objects. A projective $P$ is a retract of the coproduct of subterminals $\coprod_{u,\,s\in P(u)}y(u)$, through a section of the canonical epimorphism; if $P$ is connected the section lands in one summand, so $P$ is a retract of a subterminal, hence subterminal, $P=y(p)$. Then $p$ is completely join-prime: if $p\le\bigvee S$, the epimorphism $\coprod_{s\in S}y(s\wedge p)\to y(p)$ splits, and by connectedness $p\le s$ for one $s$. Every subterminal is a join of representables, so $L$ is join-generated by completely join-prime elements. The converse is the standard fact that sheaves on the frame of down-sets of a poset $Q$ are the presheaves on $Q$; it is not used below.
\end{proof}

So every completely distributive frame that is not of this form gives a counterexample to the question of Problem~\ref{p:presheaf}.

\subsection{The counterexample}\label{app:counterexample}

Let $X=[0,1)$ with the open sets $[0,a)$, $0\le a\le1$. Its frame is $O(X)\cong[0,1]$, a complete chain. (The same works for $\mathbf R$ with the open sets $(a,\infty)$, whose frame is the chain $[-\infty,\infty]$.)

\emph{Sheaves.} A sheaf on $X$ is a functor $F\colon((0,1],\le)^{\mathrm{op}}\to\Set$ with $F(a)\cong\lim_{0<b<a}F(b)$ for every $a$, together with $F(0)=1$: a cover of $[0,a)$ not containing $[0,a)$ is cofinal in $\{[0,b):b<a\}$, and the intersection of two members of a chain is the smaller one. Products are pointwise. Coproducts are pointwise on $(0,1]$: a compatible family over a cofinal chain in a pointwise coproduct stays in a single summand, since restriction preserves the summand index, so the pointwise coproduct, with value $1$ at $0$, is already a sheaf.

\emph{(B) holds.} At $a>0$ both sides of (B) are $\coprod_\varphi\prod_iX_{i\varphi(i)}(a)$, by distributivity in $\Set$; at $0$ both are $1$. The cases where some $J_i$ or $I$ is empty are covered, since $0$ is strict and empty products are $1$. Equivalently, every nonzero element of a chain is connected ($b\wedge c=0$ in a chain forces $b=0$ or $c=0$), so $\mathrm{Sh}(X)$ is locally connected and Proposition~\ref{a:B} applies.

\emph{(A) holds.} A complete chain is completely distributive: given $d<\bigwedge_i\bigvee_ka_{ik}$, choose for each $i$ an index $\varphi(i)$ with $a_{i\varphi(i)}>d$; then $d\le\bigwedge_ia_{i\varphi(i)}$, and the supremum over such $d$ gives the inequality $\le$. The direct argument of \S\ref{app:localic} then applies.

By Corollary~\ref{a:loc}, $\mathrm{Sh}(X)$ is also an essential subtopos of a presheaf topos, and totally distributive. Directly: $\mathrm{Sh}(X)\subseteq\mathrm{PSh}((0,1])$ is the subcategory of the functors continuous from below; the associated sheaf, $aF(a)=\lim_{b<a}F(b)$, is itself a limit, so it visibly preserves limits. In Kelly and Lawvere's terms the idempotent ideal is that of the strict inequalities $b<a$ \cite{kl}, \cite[1.5.6]{lw}.

\begin{proposition}\label{a:notpsh}
$\mathrm{Sh}(X)$ is not equivalent to a presheaf topos. It has no projective object other than $0$.
\end{proposition}

\begin{proof}
First argument: the frame $[0,1]$ has no completely join-prime element, since $0$ is the empty join and $x=\sup[0,x)$ for $x>0$. Corollary~\ref{a:loc} concludes.

Second argument. Let $P$ be projective, and let
\[\textstyle e\colon\coprod_{(c,\,z\in P(c),\,c'<c)}y(c')\longrightarrow P\]
be given by $z|_{c'}$ on the summand $(c,z,c')$. It is an epimorphism, because $\sup\{c':c'<c\}=c$. Choose a section $\sigma$ of $e$, and suppose $x\in P(v)$ with $v>0$. Since $v$ is connected, $\sigma(x)$ lies in a single summand $(c,z,c')$, with $v\le c'<c$, and $e\sigma=\mathrm{id}$ gives $x=z|_v$. The section $\sigma(z)$ over $c$ also lies in a single summand $(c_1,z_1,c_1')$, with $c\le c_1'$, and so does its restriction $\sigma(z)|_v=\sigma(z|_v)=\sigma(x)$. Hence $(c_1,z_1,c_1')=(c,z,c')$ and $c\le c'<c$, a contradiction. So $P(v)=\varnothing$ for all $v>0$, that is, $P=0$.
\end{proof}

Every nondegenerate presheaf topos has nonzero projectives, the representables, so either argument shows that the answer to Problem~\ref{p:presheaf} is no.

An object is \emph{connected} if it is not initial and is not a coproduct of two non-initial objects. For connected $Q$, the functor $\operatorname{Hom}(Q,-)$ preserves coproducts, as shown in the proof of Proposition~\ref{a:B}.

\begin{proposition}\label{a:proj}
For a Grothendieck topos $E$ the following are equivalent: (1) $E$ is equivalent to a presheaf topos $\widehat C$, $C$ small; (2) $E$ has a generating set of connected projective objects; (3) $E$ satisfies (B) and has enough projectives. In particular the presheaf topoi are the topoi satisfying (A) and (B) with enough projectives, (A) being then automatic.
\end{proposition}

\begin{proof}
(1)$\Rightarrow$(3): (B) holds by Proposition~\ref{a:ess}, applied to the identity of $\widehat C$. A representable $yc$ is projective, since $\operatorname{Hom}(yc,-)$ is evaluation at $c$ and epimorphisms of $\widehat C$ are pointwise; so every $X$ is covered by the projective $\coprod_{c,\,x\in X(c)}yc$.

(3)$\Rightarrow$(2): by Proposition~\ref{a:B}, $E$ is locally connected, $\pi_0\dashv\Delta$. Let $P$ be projective, $\eta\colon P\to\Delta\pi_0P=\coprod_{\pi_0P}1$ the unit, and $P_k$ its fibre over $k\in\pi_0P$, so that $P=\coprod_kP_k$. Decompositions of an object $X$ into two summands are the maps $X\to1\sqcup1=\Delta2$, and $\operatorname{Hom}(P,\Delta2)\cong\operatorname{Hom}(\pi_0P,2)$ through $\chi\mapsto\Delta\chi\circ\eta$; so every summand of $P$ is $\coprod_{k\in V}P_k$ for some $V\subseteq\pi_0P$. A summand $Q$ of $P_k$ is a summand of $P$, hence of this form, and $P_l\subseteq Q\subseteq P_k$ forces $P_l=0$ for $l\ne k$; so $Q$ is $0$ or $P_k$, and each nonzero $P_k$ is connected. Each $P_k$ is projective: given an epimorphism $e\colon Y\to Z$ and $q\colon P_k\to Z$, write $P=P_k\sqcup P'$ and lift $q\sqcup\mathrm{id}\colon P\to Z\sqcup P'$ along the epimorphism $e\sqcup\mathrm{id}$ to $l\colon P\to Y\sqcup P'$; the part of $P_k$ that $l$ sends into $P'$ maps to $Z\times_{Z\sqcup P'}P'=0$, so it is $0$, and $l|_{P_k}$ lifts $q$ through $e$. Now choose a generating set $G$, and for each $g\in G$ an epimorphism $P_g\to g$ with $P_g$ projective; the nonzero components of the $P_g$ form a generating set of connected projectives.

(2)$\Rightarrow$(1): let $C\subseteq E$ be the full subcategory on such a generating set. By the comparison lemma, $E\simeq\mathrm{Sh}(C,J)$, where a sieve $S$ on $c$ covers if and only if $\coprod_{(d,s)\in S}d\to c$ is an epimorphism. If it is, it has a section, since $c$ is projective; the section factors through one summand $(d,s)$, since $c$ is connected, so $s\sigma=\mathrm{id}_c$ for some $\sigma\colon c\to d$, and $\mathrm{id}_c\in S$ because sieves are closed under composition. So only maximal sieves cover, and $\mathrm{Sh}(C,J)=\widehat C$. The last sentence follows with Theorem~\ref{a:main}.
\end{proof}

The example is locally connected and satisfies (A) and (B), so it has the infinitary exactness properties of $\Set$ that he lists. What fails is the projectivity of $1$: $\Gamma=\operatorname{Hom}(1,-)$ does not preserve epimorphisms, which is not a commutation of limits with colimits. Whether $\mathrm{Sh}(X)$ has \emph{every} exactness property of $\Set$ in his informal sense is not settled here, the notion not being defined precisely.

\subsection{Independence of (A) and (B)}

\emph{(A) without (B): $\mathrm{Sh}(K)$, $K$ a Stone space such as the Cantor set.} For (A), every clopen $v$ has, for any cover, a finite clopen refinement, which can be made disjoint; so the criterion of \S\ref{app:localic} holds with clopen elements $v$. Condition (B) fails because $\mathrm{Sh}(K)$ is not locally connected; explicitly, $\prod_{n}\Delta2$ has the identity $K\to2^{\mathbf N}$ as a global section, while the global sections of $\coprod_{2^{\mathbf N}}1$ are the locally constant maps.

\emph{(B) without (A): $\mathrm{Sh}(\mathbf R)$.} It is locally connected, so (B) holds. Let $U_{n,k}=(k/2n,\,k/2n+1/n)$, $k\in\mathbf Z$. Then $f_n\colon\coprod_ky(U_{n,k})\to1$ is an epimorphism for every $n$. A section of $\prod_n\coprod_ky(U_{n,k})$ over a connected open interval $w$ forces $w\subseteq U_{n,k}$ for some $k$, for every $n$, so $w$ is empty. Hence $\prod_nf_n$ has empty domain, and it is not an epimorphism.

\emph{(B) without (A), not localic: continuous $\mathbf Z_p$-sets.} Orbits are connected, so (B) holds. Let $f_n\colon\mathbf Z_p/p^n\mathbf Z_p\to1$. The product of the $f_n$ is computed as the smooth part of the product of sets; its domain has no smooth points, since the stabiliser of every point is $\bigcap_np^n\mathbf Z_p=0$, which is not open. So the domain is empty, and $\prod_nf_n$ is not an epimorphism.

\medskip
What is not addressed here: constructive versions without the axiom of choice; whether (B) with an abstract, non-canonical isomorphism suffices; and the precise meaning of \q{all exactness properties of $\Set$}.

\begin{thebibliography}{99}
\bibitem{adelman} M.~Adelman, Abelian categories over additive ones, \emph{J.\ Pure Appl.\ Algebra} 3 (1973), 103--117, \href{https://doi.org/10.1016/0022-4049(73)90026-1}{doi:10.1016/0022-4049(73)90026-1}.
\bibitem{ar} J.~Adámek, J.~Rosický, \emph{Locally Presentable and Accessible Categories}, LMS Lecture Note Ser.~189, Cambridge University Press, 1994, \href{https://doi.org/10.1017/CBO9780511600579}{doi:10.1017/CBO9780511600579}.
\bibitem{art} J.~Adámek, J.~Rosický, V.~Trnková, Are all limit-closed subcategories of locally presentable categories reflective?, in \emph{Categorical Algebra and its Applications}, LNM~1348, Springer, 1988, 1--18, \href{https://doi.org/10.1007/BFb0081345}{doi:10.1007/BFb0081345}.
\bibitem{bkp} R.~Blackwell, G.~M.~Kelly, A.~J.~Power, Two-dimensional monad theory, \emph{J.\ Pure Appl.\ Algebra} 59 (1989), 1--41, \href{https://doi.org/10.1016/0022-4049(89)90160-6}{doi:10.1016/0022-4049(89)90160-6}.
\bibitem{caramello} O.~Caramello, Lattices of theories, \href{https://arxiv.org/abs/0905.0299}{arXiv:0905.0299}; \emph{Theories, Sites, Toposes}, Oxford University Press, 2018.
\bibitem{ccm} A.~Carboni, R.~Celia Magno, The free exact category on a left exact one, \emph{J.\ Austral.\ Math.\ Soc.\ Ser.~A} 33 (1982), 295--301, \href{https://doi.org/10.1017/S1446788700018735}{doi:10.1017/S1446788700018735}.
\bibitem{venken} J.~Venken, Non effectivité de la descente de modules plats par un morphisme fini d'anneaux locaux artiniens, \emph{C.\ R.\ Acad.\ Sci.\ Paris Sér.~A} 272 (1971), A1553--A1554.
\bibitem{three} M.~Hua, Grothendieck's \q{problem three}: generation under limits and algebraic theories relative to a class of limits, 2026, \url{https://grothendieck.commutator.io/article/grothendieck-problem-three.pdf}.
\bibitem{effective} M.~Hua, Non-effective finite equivalence relations on finite schemes over a field: Venken's example checked, and a second one, 2026, \url{https://grothendieck.commutator.io/article/grothendieck-effective.pdf}.
\bibitem{qcoh} M.~Hua, Descent of quasi-coherent sheaves along universal effective epimorphisms of schemes, 2026, \url{https://grothendieck.commutator.io/article/grothendieck-qcoh-descent.pdf}.
\bibitem{fw} B.~Fawcett, R.~J.~Wood, Constructive complete distributivity.~I, \emph{Math.\ Proc.\ Cambridge Philos.\ Soc.} 107 (1990), 81--89, \href{https://doi.org/10.1017/S0305004100068377}{doi:10.1017/S0305004100068377}.
\bibitem{freyd} P.~Freyd, Representations in abelian categories, in \emph{Proc.\ Conf.\ Categorical Algebra (La Jolla, 1965)}, Springer, 1966, 95--120, \href{https://doi.org/10.1007/978-3-642-99902-4_4}{doi:10.1007/978-3-642-99902-4\_4}.
\bibitem{gu} P.~Gabriel, F.~Ulmer, \emph{Lokal präsentierbare Kategorien}, LNM~221, Springer, 1971, \href{https://doi.org/10.1007/BFb0059396}{doi:10.1007/BFb0059396}.
\bibitem{giraud} J.~Giraud, \emph{Cohomologie non abélienne}, Grundlehren~179, Springer, 1971, \href{https://doi.org/10.1007/978-3-662-62103-5}{doi:10.1007/978-3-662-62103-5}.
\bibitem{hakim} M.~Hakim, \emph{Topos annelés et schémas relatifs}, Ergebnisse der Math.~64, Springer, 1972, \href{https://doi.org/10.1007/978-3-662-59155-0}{doi:10.1007/978-3-662-59155-0}.
\bibitem{jacqmin} P.-A.~Jacqmin, Z.~Janelidze, On stability of exactness properties under the pro-completion, \emph{Adv.\ Math.} 377 (2021), 107484, \href{https://doi.org/10.1016/j.aim.2020.107484}{doi:10.1016/j.aim.2020.107484}; \href{https://arxiv.org/abs/2002.02204}{arXiv:2002.02204}.
\bibitem{jm} P.~T.~Johnstone, I.~Moerdijk, Local maps of toposes, \emph{Proc.\ London Math.\ Soc.~(3)} 58 (1989), 281--305, \href{https://doi.org/10.1112/plms/s3-58.2.281}{doi:10.1112/plms/s3-58.2.281}.
\bibitem{stone} P.~T.~Johnstone, \emph{Stone Spaces}, Cambridge Stud.\ Adv.\ Math.~3, Cambridge University Press, 1982.
\bibitem{kellybook} G.~M.~Kelly, \emph{Basic Concepts of Enriched Category Theory}, LMS Lecture Note Ser.~64, Cambridge University Press, 1982; reprinted in \emph{Repr.\ Theory Appl.\ Categ.} 10 (2005).
\bibitem{kelly} G.~M.~Kelly, A unified treatment of transfinite constructions for free algebras, free monoids, colimits, associated sheaves, and so on, \emph{Bull.\ Austral.\ Math.\ Soc.} 22 (1980), 1--83, \href{https://doi.org/10.1017/S0004972700006353}{doi:10.1017/S0004972700006353}.
\bibitem{kl} G.~M.~Kelly, F.~W.~Lawvere, On the complete lattice of essential localizations, \emph{Bull.\ Soc.\ Math.\ Belg.\ Sér.~A} 41 (1989), no.~2, 289--319, zbMATH 0686.18005.
\bibitem{lw} R.~B.~B.~Lucyshyn-Wright, Totally distributive toposes, \emph{J.\ Pure Appl.\ Algebra} 216 (2012), 2425--2431, \href{https://doi.org/10.1016/j.jpaa.2012.03.018}{doi:10.1016/j.jpaa.2012.03.018}; \href{https://arxiv.org/abs/1108.4032}{arXiv:1108.4032}.
\bibitem{mp} M.~Makkai, R.~Paré, \emph{Accessible Categories: The Foundations of Categorical Model Theory}, Contemp.\ Math.~104, AMS, 1989, \href{https://doi.org/10.1090/conm/104}{doi:10.1090/conm/104}.
\bibitem{mr} M.~Makkai, G.~E.~Reyes, \emph{First Order Categorical Logic}, LNM~611, Springer, 1977, \href{https://doi.org/10.1007/BFb0066201}{doi:10.1007/BFb0066201}.
\bibitem{mesablishvili2004} B.~Mesablishvili, More on descent theory for schemes, \emph{Georgian Math.\ J.} 11 (2004), 783--800.
\bibitem{mesablishvili-acs} B.~Mesablishvili, Descent theory for schemes, \emph{Appl.\ Categ.\ Structures} 12 (2004), 485--512, \href{https://doi.org/10.1023/B:APCS.0000049314.33172.0d}{doi:10.1023/B:APCS.0000049314.33172.0d}.
\bibitem{mesablishvili} B.~Mesablishvili, Pure morphisms of commutative rings are effective descent morphisms for modules---a new proof, \emph{Theory Appl.\ Categ.} 7 (2000), no.~3, 38--42, \url{http://www.tac.mta.ca/tac/volumes/7/n3/7-03abs.html}.
\bibitem{moerdijk} I.~Moerdijk, \emph{Classifying Spaces and Classifying Topoi}, LNM~1616, Springer, 1995, \href{https://doi.org/10.1007/BFb0094441}{doi:10.1007/BFb0094441}.
\bibitem{ohm} J.~Ohm, R.~L.~Pendleton, Rings with noetherian spectrum, \emph{Duke Math.\ J.} 35 (1968), 631--639, \href{https://doi.org/10.1215/S0012-7094-68-03565-5}{doi:10.1215/S0012-7094-68-03565-5}.
\bibitem{pitz} M.~F.~Pitz, R.~Suabedissen, Compactifications of $\omega^*\setminus\{x\}$ and $S_\kappa\setminus\{x\}$, \emph{Topology Proc.} 45 (2015), 151--173, zbMATH 1315.54023; \href{https://arxiv.org/abs/1310.0678}{arXiv:1310.0678} (numbering of the arXiv version).
\bibitem{pr} R.~Paré, J.~Rosický, Colimits of accessible categories, \emph{Math.\ Proc.\ Cambridge Philos.\ Soc.} 155 (2013), 47--50, \href{https://doi.org/10.1017/S0305004113000030}{doi:10.1017/S0305004113000030}; \href{https://arxiv.org/abs/1110.0767}{arXiv:1110.0767}.
\bibitem{perrin} D.~Perrin, Approximation des schémas en groupes, quasi compacts sur un corps, \emph{Bull.\ Soc.\ Math.\ France} 104 (1976), 323--335, \href{https://doi.org/10.24033/bsmf.1830}{doi:10.24033/bsmf.1830}.
\bibitem{repo} Commutator-IO, grothendieck-project, \url{https://github.com/Commutator-IO/grothendieck-project}: index of results of the three courses (\texttt{docs/buffalo-*/index.tex}) and the Buffalo tab of \url{https://grothendieck.commutator.io}.
\bibitem{rw} R.~Rosebrugh, R.~J.~Wood, An adjoint characterization of the category of sets, \emph{Proc.\ Amer.\ Math.\ Soc.} 122 (1994), 409--413, \href{https://doi.org/10.1090/S0002-9939-1994-1216823-2}{doi:10.1090/S0002-9939-1994-1216823-2}.
\bibitem{rydh} D.~Rydh, Submersions and effective descent of étale morphisms, \emph{Bull.\ Soc.\ Math.\ France} 138 (2010), 181--230, \href{https://doi.org/10.24033/bsmf.2588}{doi:10.24033/bsmf.2588}.
\bibitem{stacks} The Stacks Project Authors, \emph{The Stacks Project}, \url{https://stacks.math.columbia.edu}, tags as cited.
\bibitem{wilson} T.~M.~Wilson, Weak Vop\v{e}nka's principle does not imply Vop\v{e}nka's principle, \emph{Adv.\ Math.} 363 (2020), 106986, \href{https://doi.org/10.1016/j.aim.2020.106986}{doi:10.1016/j.aim.2020.106986}.
\end{thebibliography}
\end{document}
